\documentclass[12pt]{article}
\usepackage[utf8]{inputenc}
\usepackage[T1]{fontenc}
\usepackage{amsmath, amssymb, amsthm}

\title{Teorema Fundamental do Cálculo}
\author{Notas de Estudo}
\date{}

% Ambientes de teorema
\newtheorem{theorem}{Teorema}
\newtheorem{lemma}[theorem]{Lema}
\newtheorem{corollary}[theorem]{Corolário}
\newtheorem{definition}{Definição}
\newtheorem{example}{Exemplo}

\begin{document}
\maketitle

\section{Introdução}

O Teorema Fundamental do Cálculo estabelece a conexão profunda entre
diferenciação e integração — duas operações aparentemente distintas.

\section{Enunciados}

\begin{theorem}[TFC — Parte 1]
  Seja $f$ contínua em $[a,b]$. A função:
  \[
  F(x) = \int_a^x f(t)\,dt
  \]
  é diferenciável em $(a,b)$ e $F'(x) = f(x)$.
\end{theorem}

\begin{proof}
  Pela definição de derivada:
  \[
  F'(x) = \lim_{h \to 0} \frac{F(x+h) - F(x)}{h}
         = \lim_{h \to 0} \frac{1}{h}\int_x^{x+h} f(t)\,dt
  \]
  Pelo Teorema do Valor Médio para integrais, existe $c \in (x, x+h)$ tal que:
  \[
  \frac{1}{h}\int_x^{x+h} f(t)\,dt = f(c)
  \]
  Como $c \to x$ quando $h \to 0$ e $f$ é contínua, $f(c) \to f(x)$.
  Portanto $F'(x) = f(x)$. \qed
\end{proof}

\begin{theorem}[TFC — Parte 2]
  Se $F$ é uma primitiva de $f$ em $[a,b]$, então:
  \[
  \int_a^b f(x)\,dx = F(b) - F(a)
  \]
\end{theorem}

\section{Exemplos}

\begin{example}
  Calcule $\displaystyle\int_0^\pi \sin x\,dx$.

  Como $F(x) = -\cos x$ é primitiva de $\sin x$:
  \[
  \int_0^\pi \sin x\,dx = -\cos(\pi) - (-\cos 0) = 1 + 1 = 2
  \]
\end{example}

\begin{example}
  A derivada de $\displaystyle G(x) = \int_1^x e^{t^2}\,dt$ é:
  \[
  G'(x) = e^{x^2}
  \]
  diretamente pela Parte 1 do TFC, sem precisar calcular a integral.
\end{example}

\section{Corolários}

\begin{corollary}
  Toda função contínua possui primitiva.
\end{corollary}

\begin{corollary}
  Se $f$ é contínua e $\displaystyle\int_a^b f(x)\,dx = 0$ para todo $[a,b]$,
  então $f \equiv 0$.
\end{corollary}

\end{document}
